\uuid{pQZK}
\exo7id{2726}
\titre{exo7 2726}
\auteur{matexo1}
\organisation{exo7}
\datecreate{2002-02-01}
\isIndication{false}
\isCorrection{false}
\chapitre{Série numérique}
\sousChapitre{Autre}
\cpge{mpsi-mp2i.procedes-sommatoires-discrets.a, pcsi.series-numeriques.a, ptsi.series-numeriques.a, mpsi-mp2i.procedes-sommatoires-discrets.b, pcsi.series-numeriques.b, ptsi.series-numeriques.b, mpsi-mp2i.procedes-sommatoires-discrets.c, pcsi.series-numeriques.c, ptsi.series-numeriques.c, bcpst1.derivees-primitives-integrales.d, mp-mpi.series-numeriques-vectorielles.a, pt.series-numeriques.b, mp-mpi.series-numeriques-vectorielles.b, pc.series-numeriques, psi.series-numeriques.a, pt.series-numeriques.a, bcpst2.series-reelles}
\module{Analyse}
\niveau{L1}
\difficulte{}

\contenu{
\question{
En comparant
avec les int\'egrales de Wallis
$$I_n = \int_0^{\pi /2} \cos^n t\,dt,$$
d\'eterminer la nature et calculer
la somme de la s\'erie
$$\displaystyle \sum_{n\geq 1} (-1)^n 2^{-2n} C_{2n}^n.$$
}
}
